% =============================================================================
% 1. 基础设置
% =============================================================================
\documentclass[twocolumn, a4paper, 10pt, UTF8]{ctexart}

% --- 页面布局 ---
\usepackage{geometry}
\geometry{left=2.5cm, right=2.5cm, top=2.5cm, bottom=2.5cm}
\setlength{\columnsep}{2em} % 两栏间距

% --- 清除页眉页脚 ---
\pagestyle{empty}

% --- 字体与数学包 ---
\usepackage{times}      % Times New Roman 风格
\usepackage{mathptmx}   % 数学公式字体
\usepackage{amsmath, amsfonts, amssymb, bm}

% --- 图表与排版 ---
\usepackage{graphicx}
\usepackage{booktabs}   % 三线表
\usepackage{caption}
\usepackage{bicaption}  % 双语标题
\usepackage{stfloats}   % 双栏底部浮动支持
\usepackage{multirow}   % 多行合并
\usepackage{pgfplots}   % 绘图包
\usepackage{tikz}       % 绘图包
\pgfplotsset{compat=1.18}

% --- 标题格式定制 ---
\usepackage{titlesec}
\titleformat{\section}{\large\heiti\bfseries}{\thesection}{1em}{}
\titleformat{\subsection}{\normalsize\heiti\bfseries}{\thesubsection}{1em}{}
\titleformat{\subsubsection}{\small\heiti\bfseries}{\thesubsubsection}{1em}{}

% --- 标题间距调整 ---
\titlespacing*{\section}{0pt}{1.5ex plus 0.5ex minus 0.2ex}{1ex plus 0.2ex}
\titlespacing*{\subsection}{0pt}{1ex plus 0.3ex minus 0.1ex}{0.8ex plus 0.2ex}
\titlespacing*{\subsubsection}{0pt}{0.8ex plus 0.2ex minus 0.1ex}{0.6ex plus 0.1ex}

% --- 图表标题设置 ---
\captionsetup{font={small,bf}, labelfont=bf, labelsep=quad, justification=centering}
\captionsetup[figure][bi-first]{name=图}
\captionsetup[figure][bi-second]{name=Fig.}
\captionsetup[table][bi-first]{name=表}
\captionsetup[table][bi-second]{name=Table}

% --- 段落缩进设置 ---
\setlength{\parindent}{2em}  % 设置段落首行缩进为2个字符
\setlength{\parskip}{0pt}    % 段落间距设为0

% 强制启用段落缩进（针对ctex文档类）
\usepackage{indentfirst}

% --- 辅助命令 ---
\newcommand{\txt}[1]{\mathrm{#1}}
\newcommand{\unitdot}{\cdot}

\begin{document}

% =============================================================================
% 2. 通栏区域 (标题、摘要、短横线)
% =============================================================================
\twocolumn[
    \begin{@twocolumnfalse}
        \thispagestyle{empty}

        % --- 中文头部 ---
        \begin{center}
            \vspace{1em}
            {\huge \heiti \textbf{新能源取代化石燃料成为能源主流的时间进程与条件分析}} \par
            \vspace{1.5em}
            {\large \fangsong 刘培宽$^{1,a}$} \par
            \vspace{0.5em}
            {\footnotesize \songti (1. 上海工程技术大学 a. 电子电气工程学院)}
        \end{center}

        \vspace{1em}

        % --- 中文摘要 ---
        \noindent {\small \textbf{摘要：}在全球减排目标和能源转型持续推进的背景下，风电、光伏等新能源的发展速度明显加快，“新能源是否以及何时能够取代化石燃料成为能源主流”也成为一个被频繁讨论的问题。不过，从现实情况来看，这一问题并不存在一个放之四海而皆准的答案，其结论往往取决于具体的分析层面和时间尺度。基于这一认识，本文结合国际能源署（IEA）、国际可再生能源署（IRENA）等权威机构发布的公开统计数据和年度报告，对全球能源结构的演变情况进行梳理，并重点关注新能源发电成本变化、装机规模扩张等影响其替代能力的关键因素。在此基础上，从不同能源领域和情景假设出发，对新能源替代化石能源的时间进程进行阶段性分析。 \par
            \vspace{0.5em}
            \noindent \textbf{关键词：} 新能源；化石燃料；能源结构转型；可再生能源；度电成本（LCOE）；能源情景分析；电力结构 \par
            \vspace{0.5em}
        }

        \vspace{2em}

        % --- 英文头部 ---
        \begin{center}
            {\Large \textbf{Analysis of the Timeline and Conditions for Renewable Energy\\ to Replace Fossil Fuels as the Mainstream Energy Source}} \par
            \vspace{1em}
            {\large \textit{LIU Peikuan$^{1,a}$}} \par
            \vspace{0.5em}
            {\footnotesize \itshape (1a. College of Electronic and Electrical Engineering, Shanghai University of Engineering Science (SUES))}
        \end{center}

        % --- 英文摘要 ---
        \vspace{1em}
        \noindent {\small \textbf{Abstract:} Against the backdrop of global emission reduction targets and continuous advancement of energy transition, the development of renewable energy sources such as wind and solar power has significantly accelerated. The question of "whether and when renewable energy can replace fossil fuels as the mainstream energy source" has become a frequently discussed topic. However, from a realistic perspective, this question does not have a universally applicable answer, and its conclusion often depends on specific analytical dimensions and time scales. Based on this understanding, this paper combines publicly available statistical data and annual reports from authoritative institutions such as the International Energy Agency (IEA) and the International Renewable Energy Agency (IRENA) to analyze the evolution of global energy structure, with particular focus on key factors affecting renewable energy's substitution capacity, such as changes in renewable energy generation costs and expansion of installed capacity. On this basis, a phased analysis of the timeline for renewable energy to replace fossil fuels is conducted from different energy sectors and scenario assumptions. \par
            \vspace{0.5em}
            \noindent \textbf{Key words:} renewable energy; fossil fuels; energy structure transition; renewable energy; levelized cost of electricity (LCOE); energy scenario analysis; power structure}

        % --- 短横线 ---
        \vspace{1.5em}
        \centerline{\rule{0.5\textwidth}{0.5pt}}
        \vspace{2em}

    \end{@twocolumnfalse}
]

% =============================================================================
% 3. 正文内容 (大量填充以撑开页面)
% =============================================================================

% 确保正文段落有缩进
\setlength{\parindent}{2em}

% --- 第一章 ---
\section{全球能源结构与能源主流格局现状}

\subsection{全球能源主流格局的基本界定}
在讨论新能源是否能够取代化石燃料成为能源主流之前，有必要首先明确"能源主流"这一概念在不同分析层面上的含义。一般而言，能源主流并非单一指标，而是可以从一次能源结构、电力结构以及终端能源消费结构等多个维度进行理解。在较长时期内，煤炭、石油和天然气等化石能源在全球一次能源消费中占据主导地位，并构成现代能源系统的基础框架。随着新能源技术的发展和应用规模的扩大，这一格局正在发生变化，但不同能源形态在各个层面的演进节奏并不完全一致。

\subsection{全球电力结构中新能源与化石能源占比变化}
在全球范围内，电力结构的变化为观察能源转型趋势提供了直观依据。近年来，随着风电、太阳能光伏等新能源装机容量的持续增长，其在全球发电量中的占比逐步提升，与此同时，煤电等传统化石能源发电的相对占比呈现下降趋势。不同国家和地区由于资源禀赋、政策导向和经济发展阶段存在差异，其电力结构变化的节奏和幅度也有所不同，但整体趋势具有一定一致性。本文选取 Ember 发布的全球电力统计数据，对 2000—2024 年间不同电源类型在全球发电量中的占比变化进行对比分析，相关结果如图 1 所示。

% 图1：全球电力结构变化（2000-2024年）
\begin{figure}[htbp]
    \centering
    \begin{tikzpicture}
        \begin{axis}[
                width=0.95\linewidth,
                height=6cm,
                xlabel={年份 (20XX)},
                ylabel={占全球发电量比重(\%)},
                xmin=2000, xmax=2024,
                ymin=0, ymax=40,
                axis lines*=left,
                legend style={at={(0.5,-0.25)}, anchor=north, legend columns=3, font=\scriptsize},
                tick label style={font=\small},
                label style={font=\small},
                xtick={2000,2010,2020,2024},
                xticklabels={00,10,20,24},
                ytick={0,10,20,30,40},
                grid=major,
                grid style={gray!30},
            ]

            % 核电累积线 (实线，圆点标记) - 底层，从0累积
            \addplot[thick, solid, mark=*, mark size=2pt, black] coordinates {
                    (2000,17.0) (2010,13.0) (2020,10.0) (2024,9.0)
                };
            \addlegendentry{核电 (9\%)}

            % 核电+水电累积线 (虚线，方块标记)
            \addplot[thick, dashed, mark=square*, mark size=2pt, black] coordinates {
                    (2000,33.0) (2010,28.0) (2020,24.5) (2024,23.0)
                };
            \addlegendentry{核电+水电 (23\%)}

            % 核电+水电+其他可再生累积线 (点线，三角标记)
            \addplot[thick, dotted, mark=triangle*, mark size=2.5pt, black] coordinates {
                    (2000,35.0) (2010,31.0) (2020,27.5) (2024,26.0)
                };
            \addlegendentry{+其他可再生 (26\%)}

            % 核电+水电+其他+风能累积线 (点划线，菱形标记)
            \addplot[thick, dashdotted, mark=diamond*, mark size=2pt, black] coordinates {
                    (2000,35.1) (2010,32.8) (2020,33.5) (2024,34.0)
                };
            \addlegendentry{+风能 (34\%)}

            % 全部低碳电力累积线 (双点划线，五角星标记) - 顶层
            \addplot[thick, dashdotdotted, mark=star, mark size=2.5pt, black] coordinates {
                    (2000,35.1) (2010,32.9) (2020,36.5) (2024,41.0)
                };
            \addlegendentry{全部低碳电力 (41\%)}

        \end{axis}
    \end{tikzpicture}
    \bicaption{全球电力结构中不同电源类型占比变化（2000-2024年）}
    {Changes in the proportion of different power sources in global electricity structure (2000-2024)}
    \label{fig:power-structure}
\end{figure}

\subsection{不同能源形态下"主流地位"的差异}
需要指出的是，新能源在电力领域占比的提升，并不等同于其在整个能源系统中已经全面取代化石能源成为主流。一次能源消费结构仍然受到化石能源存量体系、基础设施以及终端用能方式等多方面因素的影响。在部分发展中经济体和资源型地区，化石能源在相当长一段时间内仍将发挥重要作用。因此，在评价新能源的"主流地位"时，有必要区分分析层面，避免将电力结构的变化简单外推至整个能源系统。

% --- 第二章 ---
\section{新能源替代化石能源的客观驱动因素}
\subsection{新能源发电成本变化趋势}
根据国际可再生能源署（IRENA）发布的《Renewable Power Generation Costs in 2024》报告，2010—2024 年期间，全球主要可再生能源技术的平准化度电成本（Levelized Cost of Electricity，LCOE）整体呈现显著下降趋势。其中，太阳能光伏和陆上风电的成本降幅尤为明显，已由早期明显高于传统化石能源，逐步下降至在多数地区具备成本竞争力的水平。海上风电和聚光太阳能发电等技术虽然成本下降路径相对平缓，但其长期趋势同样表现为逐步降低。成本并非逐年单调下降，但长期下降趋势明确。这些波动与原材料价格、融资环境以及供应链状况等因素密切相关。相关变化趋势如图 2 所示。

% 图2：新能源发电成本变化趋势
\begin{figure}[htbp]
    \centering
    \begin{tikzpicture}
        \begin{axis}[
                width=0.95\linewidth,
                height=6cm,
                xlabel={年份 (20XX)},
                ylabel={LCOE (2022 USD/kWh)},
                xmin=2010, xmax=2024,
                ymin=0, ymax=0.4,
                axis lines*=left,
                legend style={at={(0.5,-0.25)}, anchor=north, legend columns=2, font=\scriptsize},
                tick label style={font=\small},
                label style={font=\small},
                xtick={2010,2012,2014,2016,2018,2020,2022,2024},
                xticklabels={10,12,14,16,18,20,22,24},
                ytick={0,0.1,0.2,0.3,0.4},
                grid=major,
                grid style={gray!30},
            ]

            % 太阳能光伏 (实线，圆点标记) - 最大降幅
            \addplot[thick, solid, mark=*, mark size=2pt, black] coordinates {
                    (2010,0.40) (2011,0.35) (2012,0.25) (2013,0.19) (2014,0.18) (2015,0.13)
                    (2016,0.10) (2017,0.08) (2018,0.07) (2019,0.06) (2020,0.05) (2021,0.05)
                    (2022,0.04) (2023,0.03) (2024,0.03)
                };
            \addlegendentry{太阳能光伏 (PV)}

            % 海上风电 (虚线，方块标记) - 中等降幅
            \addplot[thick, dashed, mark=square*, mark size=2pt, black] coordinates {
                    (2010,0.20) (2011,0.21) (2012,0.18) (2013,0.15) (2014,0.18) (2015,0.15)
                    (2016,0.12) (2017,0.11) (2018,0.10) (2019,0.08) (2020,0.08) (2021,0.07)
                    (2022,0.07) (2023,0.06) (2024,0.07)
                };
            \addlegendentry{海上风电}

            % 陆上风电 (点线，三角标记) - 稳定降幅
            \addplot[thick, dotted, mark=triangle*, mark size=2.5pt, black] coordinates {
                    (2010,0.10) (2011,0.09) (2012,0.09) (2013,0.08) (2014,0.08) (2015,0.07)
                    (2016,0.06) (2017,0.05) (2018,0.04) (2019,0.04) (2020,0.03) (2021,0.03)
                    (2022,0.03) (2023,0.03) (2024,0.03)
                };
            \addlegendentry{陆上风电}

            % 聚光太阳能 (点划线，菱形标记) - 波动较大
            \addplot[thick, dashdotted, mark=diamond*, mark size=2pt, black] coordinates {
                    (2010,0.39) (2011,0.37) (2012,0.36) (2013,0.27) (2014,0.23) (2015,0.24)
                    (2016,0.27) (2017,0.22) (2018,0.16) (2019,0.23) (2020,0.12) (2021,0.12)
                    (2022,0.12) (2023,0.11) (2024,0.08)
                };
            \addlegendentry{聚光太阳能 (CSP)}

        \end{axis}
    \end{tikzpicture}
    \bicaption{全球主要可再生能源技术LCOE变化趋势（2010-2024年）}
    {LCOE trends of major renewable energy technologies globally (2010-2024)}
    \label{fig:lcoe-trends}
\end{figure}

\subsection{新能源装机规模与发电量增长特征}
IRENA 的统计数据显示，过去二十年间，全球可再生能源新增装机容量保持持续增长态势，且增长速度明显加快。尤其是在近几年，新能源已成为全球新增发电装机的主要来源。从结构上看，太阳能光伏和风电是新增装机的主要贡献者。这一方面与其成本优势密切相关，另一方面也反映出其技术成熟度和部署灵活性不断提升。尽管装机容量与实际发电量之间并非完全线性关系，但随着设备效率提高和运行经验积累，新能源在全球电力供给中的实际贡献持续上升。相关装机增长特征如图 3 所示。

\begin{figure}[htbp]
    \centering
    \begin{tikzpicture}
        \begin{axis}[
                width=0.95\linewidth,
                height=6cm, % ✅ 跟图1一致：给底部空间
                xlabel={年份 (20XX)},
                xlabel style={font=\small}, % ✅ 删掉 yshift，别再往下挤
                ylabel={新增装机容量 (GW/年)},
                xmin=2013.5, xmax=2024.5,
                ymin=0, ymax=620,
                axis lines*=left,
                legend style={at={(0.5,-0.25)}, anchor=north, legend columns=1, font=\scriptsize},
                tick label style={font=\small},
                label style={font=\small},
                xtick={2014,2015,2016,2017,2018,2019,2020,2021,2022,2023,2024},
                xticklabels={14,15,16,17,18,19,20,21,22,23,24},
                ytick={0,100,200,300,400,500,600},
                grid=major,
                grid style={gray!30},
                ybar,
                bar width=0.6,
            ]
            \addplot[fill=gray!50, draw=black, thick] coordinates {
                    (2014,127) (2015,153) (2016,169) (2017,168) (2018,171)
                    (2019,185) (2020,270) (2021,263) (2022,303) (2023,484) (2024,582)
                };
            \addlegendentry{新能源新增装机容量}

            % 添加数值标签
            \node at (axis cs:2014,127) [above] {\scriptsize 127};
            \node at (axis cs:2015,153) [above] {\scriptsize 153};
            \node at (axis cs:2016,169) [above] {\scriptsize 169};
            \node at (axis cs:2017,168) [above] {\scriptsize 168};
            \node at (axis cs:2018,171) [above] {\scriptsize 171};
            \node at (axis cs:2019,185) [above] {\scriptsize 185};
            \node at (axis cs:2020,270) [above] {\scriptsize 270};
            \node at (axis cs:2021,263) [above] {\scriptsize 263};
            \node at (axis cs:2022,303) [above] {\scriptsize 303};
            \node at (axis cs:2023,484) [above] {\scriptsize 484};
            \node at (axis cs:2024,582) [above] {\scriptsize 582};
        \end{axis}
    \end{tikzpicture}
    \bicaption {全球可再生能源新增装机容量变化趋势}
    {Global renewable energy capacity additions trend}
    \label{fig:capacity-growth}
\end{figure}


\subsection{技术进步与产业规模效应对能源结构的影响}
新能源替代能力的形成不仅体现在成本和规模层面，还与技术进步及产业成熟度的提升密切相关。从技术角度看，光伏组件和风电设备在转换效率、可靠性以及运维水平等方面持续改进，使单位装机容量所能提供的有效电量不断提高。从产业层面看，新能源产业链的规模化发展和全球化布局，使设备制造成本、建设周期和运维成本持续下降，并形成明显的规模经济效应。然而，IRENA 的成本分解分析也表明，新能源发电成本并不完全由设备成本决定，融资条件在不同地区对 LCOE 的影响尤为显著。这一现象说明，尽管新能源技术本身已具备较强的成本竞争力，但其在不同地区、不同能源系统中的替代能力仍受到金融环境、制度条件和基础设施水平等多重因素制约。因此，技术进步和产业规模效应为新能源替代化石能源提供了重要支撑，但并不意味着这一过程在全球范围内可以同步、无差异地推进。相关成本结构差异如表 1 所示。

% 表1：新能源发电成本结构差异
\begin{table}[htbp]
    \centering
    \bicaption{不同地区新能源发电成本中融资成本占比对比}
    {Comparison of financing cost share in renewable energy LCOE across different regions}
    \label{tab:cost-structure}
    \small
    \begin{tabular}{lccc}
        \toprule
        \multirow{2}{*}{\textbf{地区}} & \multicolumn{2}{c}{\textbf{融资成本占比}} & \multirow{2}{*}{\textbf{平均占比}}               \\
        \cmidrule(lr){2-3}
                                     & \textbf{陆上风电/\%}                    & \textbf{太阳能光伏/\%}              &             \\
        \midrule
        亚洲                           & 18                                  & 25                             & 22          \\
        欧洲                           & 21                                  & 31                             & 26          \\
        北美                           & 35                                  & 42                             & 39          \\
        南美                           & 28                                  & 33                             & 31          \\
        非洲                           & 45                                  & 48                             & 47          \\
        大洋洲                          & 24                                  & 29                             & 27          \\
        \midrule
        \textbf{全球平均}                & \textbf{29}                         & \textbf{33}                    & \textbf{31} \\
        \bottomrule
    \end{tabular}

    \vspace{0.5em}
    \fbox{\parbox{0.9\linewidth}{\small
            \textbf{融资成本的隐性影响}\\
            被认为"高风险"的国家面临更高的资本成本，推高了融资成本在平准化度电成本(LCOE)中的占比。数据显示，融资条件对不同地区新能源项目的经济性产生显著差异化影响。
        }}
\end{table}

% --- 第三章 ---
\section{新能源取代化石能源的时间进程分析}
\subsection{不同能源领域中新能源替代节奏的差异}
如果从整个能源系统内部来看，不同领域中新旧能源更替的节奏差异是相当明显的。电力领域往往被认为是新能源推进最快的部分，这并不难理解。一方面，电力系统的能源形态相对集中，技术替代路径也较为清晰；另一方面，光伏、风电等技术在成本持续下降的同时，装机规模快速扩张，使得新能源在新增电力装机中逐渐占据主导地位。在一些国家和地区，这种变化已经开始反映到发电量结构之中。但如果把视角从电力系统拉回到更广义的一次能源消费结构，情况就要复杂得多了。化石能源目前仍然在全球一次能源消费中占据相当大的比重，这既与交通、工业等终端用能领域对化石能源的高度依赖有关，也与长期形成的基础设施体系密切相关。

\subsection{国际机构能源转型情景下的时间尺度判断}
针对新能源与化石能源在长期演进中的相对地位，国际能源署（IEA）、国际可再生能源署（IRENA）等机构通常采用情景分析的方法，对不同政策强度和转型条件下的能源结构变化进行系统评估。需要注意的是，这类情景分析本身并不是对未来的精确预测，而是在一系列假设前提下，对可能路径进行的理性推演。从"既有政策情景"来看，新能源在未来一段时间内仍将保持较快增长，但化石能源在一次能源消费中的主导地位难以在短期内被根本性动摇。而在更为积极的加速转型或净零排放情景中，新能源渗透速度显著提高，其在电力系统中的主导地位有可能在更早阶段形成。不过，这一过程高度依赖于政策的持续性、投资规模的稳定性，以及电网、储能等配套基础设施是否能够同步推进。

\subsection{新能源成为能源主流的阶段性特征分析}
综合以上讨论可以发现，新能源取代化石能源并成为能源主流，更像是一个由局部向整体、由增量向存量逐步推进的过程，而非一次性完成的转折点。在初期阶段，这种转型主要体现为新能源在新增装机和新增发电量中的占比不断提高，对传统化石能源发电形成边际上的替代。进入中期阶段，随着技术成熟度提升、电网适应能力增强以及相关基础设施逐步完善，新能源在电力系统中的地位将进一步巩固，并开始对化石能源发电产生更明显的结构性挤出效应。但即便如此，这一阶段仍难以在终端能源消费领域实现对化石能源的全面替代。从更长远的视角来看，新能源是否能够在整个能源系统层面真正取代化石能源，不仅取决于技术进步和成本变化，还受到金融条件、制度安排以及区域差异等多方面因素的共同影响。因此，与其将"新能源成为能源主流"理解为某一个确定年份的结果，不如将其视为一个持续推进、阶段性演化的系统性转型过程。

% --- 第四章 ---
\section{新能源成为能源主流面临的主要约束条件}
\subsection{能源系统运行与消纳能力约束}
与单一技术替代不同，能源转型首先面对的是整个能源系统的运行约束。新能源发电具有明显的波动性和间歇性特征，其出力与气象条件密切相关，这对以稳定供能为目标构建的传统电力系统提出了新的要求。即便在新能源装机容量快速增长的情况下，电力系统是否具备足够的调节能力，仍然是影响其进一步发展的关键因素。在现实运行中，电网调峰能力不足、跨区域输电受限以及储能设施建设滞后，都会在一定程度上限制新能源电量的有效消纳。一些地区即便具备良好的风能或太阳能资源条件，也可能因为系统调节能力不足而难以充分利用已有装机。这种"装机有余、消纳不足"的现象，说明新能源替代化石能源并不仅仅是增加装机规模的问题，更涉及电力系统整体结构的协同调整。

\subsection{区域资源禀赋与基础设施差异}
新能源发展在不同国家和地区之间呈现出明显的不均衡特征，这种差异不仅体现在资源禀赋上，也体现在基础设施和制度条件上。风能、太阳能等资源在空间分布上具有天然差异，一些地区具备发展新能源的自然优势，而另一些地区则受到资源条件限制，其新能源潜力相对有限。除了资源本身，基础设施条件同样发挥着重要作用。电网建设水平、跨区域输电能力以及储能设施的完善程度，都会直接影响新能源在当地能源结构中的地位。在基础设施相对薄弱的地区，即便新能源技术成本已经显著下降，其实际应用效果仍可能受到制约。

\subsection{能源转型过程中面临的不确定性因素}
除了系统性和区域性约束之外，能源转型本身还面临多种不确定性因素。这些不确定性既包括技术层面的，也涉及经济和制度层面的变化。例如，原材料价格波动、融资成本变化以及全球产业链调整，都可能在短期内影响新能源项目的成本结构和投资意愿。因此，新能源替代化石能源并成为能源主流，并不是一条线性、确定的路径，而是在多重约束和不确定性条件下不断推进、反复调整的过程。这也进一步说明，将能源转型简单理解为"某一年完成替代"，在现实分析中并不具备可行性。

% --- 第五章 ---
\section{结论}
围绕新能源是否以及在何种意义下能够取代化石能源成为能源主流，本文从全球能源结构现状、替代能力来源、时间进程以及现实约束等方面进行了系统梳理。综合前文分析可以看到，新能源的发展已经在多个层面产生了实质性影响，但这一变化并不呈现为简单、同步的"全面替代"，而是表现出明显的层级差异和阶段特征。

从能源结构的实际变化来看，电力领域仍然是新能源渗透最为迅速、结构调整最为明显的部分。风电和光伏在新增装机和发电量中的占比持续提升，使新能源逐步成为电力系统中不可忽视的重要组成。但如果将视角扩展到一次能源消费和终端用能层面，化石能源的主导地位依然稳固，新能源的替代效应尚未形成同等程度的结构性变化。

从时间进程的角度来看，新能源取代化石能源并不存在一个统一、明确的时间节点。国际机构的情景分析以及现实数据都表明，这一过程更适合理解为分领域推进、分阶段展开的长期演化。能源系统自身的运行约束、区域资源和基础设施差异，以及技术、经济和政策层面的不确定性，共同决定了新能源转型难以实现快速、线性推进。这些因素并不否定新能源发展的长期趋势，但提醒我们在评价其替代进程时，需要保持对现实条件和结构限制的充分认识。

% =============================================================================
% 4. 参考文献
% =============================================================================
\renewcommand{\refname}{\raggedright \large \heiti 参考文献：}
\begin{thebibliography}{99}
    \setlength{\itemsep}{0pt} \small
    \bibitem{1} 张恒硕. 能源转型试点政策的化石能源消耗减控效应研究[D]. 东北石油大学, 2025.
    \bibitem{2} 齐英瑛. 双碳目标下中国能源转型路径研究[D]. 四川大学, 2025.
    \bibitem{3} 刘佳琳. "双碳"背景下能源相关税收政策对能源结构转型的影响分析[D]. 东北财经大学, 2024.
    \bibitem{4} IEA. World Energy Outlook 2024[R]. Paris: IEA, 2024.
    \bibitem{5} International Renewable Energy Agency. Renewable Power Generation Costs in 2024[R]. Abu Dhabi: IRENA, 2025.
    \bibitem{6} Energy Institute. Statistical Review of World Energy 2025[R]. London: Energy Institute, 2025.
\end{thebibliography}

\end{document}